\protect \section *{\nameref *{sec:plasma-collisions}}
\begin{Solution}{4.{1}}
		$\sigma\propto T^{3/2}$: зростання в~$31{,}6$ раза; $\sigma(1~\text{кеВ})\approx3{,}6\times10^{17}~\mathrm{с}^{-1}$ (мідь: $\approx5\times10^{17}~\mathrm{с}^{-1}$); $j^2/\sigma\propto T^{-3/2}$.
	
\end{Solution}
\begin{Solution}{4.{2}}
		$\kappa_e$ зростає в~316 разів, $\sigma$ в~31,6 раза; $\kappa_e/\sigma\propto T_e$ зростає в~10 разів, а~$\kappa_e/(\sigma T_e)$ не змінюється
		(аналог закону Відемана~--- Франца).
	
\end{Solution}
\begin{Solution}{4.{3}}
		в~обох випадках $E_D\approx3{,}9$~В/м (відношення $n/T_e$ однакове);
	$v_c/v_{Te}=\sqrt{E_D/E}\approx6{,}2$, частка $\sim\exp(-E_D/2E)\approx3\times10^{-9}$.
	
\end{Solution}
\begin{Solution}{4.{4}}
		$D_a=2D_i$ при $T_e=T_i$ і~$D_a=11D_i$ при $T_e=10T_i$.
	
\end{Solution}
\begin{Solution}{4.{5}}
		іоносфера: $\ln\Lambda\approx12{,}6$, $\nu_{ei}\approx10^{3}~\mathrm{с}^{-1}$, $\nu_{ei}/\omega_{pe}\approx2\times10^{-5}$;
	корона: $\ln\Lambda\approx19$, $\nu_{ei}\approx6~\mathrm{с}^{-1}$, $\nu_{ei}/\omega_{pe}\approx10^{-8}$;
	лабораторна: $\ln\Lambda\approx11{,}5$, $\nu_{ei}\approx10^{7}~\mathrm{с}^{-1}$, $\nu_{ei}/\omega_{pe}\approx6\times10^{-5}$~--- найбільш зіткненнєва.
	
\end{Solution}
\begin{Solution}{4.{6}}
		$\tau_e\approx2\times10^{-4}$~с, $\ell_e\approx8$~км $\gg2\pi R\approx19$~м;
	$\tau_{eq}\sim0{,}2$~с для водню.
	
\end{Solution}
\begin{Solution}{4.{7}}
		$\nu_{ei}\approx5\times10^{3}~\mathrm{с}^{-1}$, $r_{Le}\approx5\times10^{-3}$~см, $D_\perp\approx0{,}1$~см$^2$/с;
	$D_B\approx1{,}3\times10^{6}$~см$^2$/с, тобто в~$\sim10^{7}$ разів більший. Якби дифузія була бомівською,
	утримання в~реакторних масштабах було б неможливим; реальний (аномальний) перенос лежить між цими межами.
	
\end{Solution}
\begin{Solution}{4.{8}}
		$\sigma=10^{7}$~См/м $=9\times10^{16}~\mathrm{с}^{-1}$, $T_e\approx0{,}4$~кеВ; $T_i\approx250$~еВ
	($\Delta\lambda/\lambda_0=1{,}22\times10^{-3}$, $T_i=(\Delta\lambda/\lambda_0)^2Mc^2/8\ln2$).
	
\end{Solution}
